\documentclass[a4paper,12pt]{article}
\usepackage[utf8]{inputenc}
\usepackage[T1]{fontenc}
\usepackage[magyar]{babel}
\usepackage{amsmath,amssymb,amsthm}
\usepackage[margin=2cm]{geometry}
\usepackage{graphicx}
\usepackage{parskip}
\usepackage{enumitem}
\usepackage{xcolor}
\usepackage{mdframed}
\usepackage{booktabs}
\usepackage{hyperref}
\hypersetup{colorlinks=true,linkcolor=blue}

\definecolor{defbg}{rgb}{0.95,1.0,0.95}
\definecolor{thmbg}{rgb}{1.0,0.96,0.92}
\definecolor{proofbg}{rgb}{0.97,0.97,1.0}
\definecolor{exbg}{rgb}{0.97,0.97,0.97}

\newtheoremstyle{mydef}{}{}{}{}{\bfseries}{.}{ }{}
\theoremstyle{mydef}
\newtheorem{tetel}{Tétel}[section]
\newtheorem{lemma}[tetel]{Lemma}
\newtheorem{definicio}[tetel]{Definíció}
\newtheorem{kovetkezmeny}[tetel]{Következmény}
\newtheorem{megjegyzes}[tetel]{Megjegyzés}
\makeatletter
\renewcommand*{\theHtetel}{\thesection.\arabic{tetel}}
\providecommand*{\theHlemma}{\theHtetel}
\providecommand*{\theHdefinicio}{\theHtetel}
\providecommand*{\theHkovetkezmeny}{\theHtetel}
\providecommand*{\theHmegjegyzes}{\theHtetel}
\makeatother


% Kényelmi makrók
\newcommand{\R}{\mathbb{R}}
\newcommand{\C}{\mathbb{C}}
\newcommand{\N}{\mathbb{N}}
\newcommand{\norm}[1]{\left\|#1\right\|}
\newcommand{\abs}[1]{\left|#1\right|}
\newcommand{\innerprod}[2]{\left\langle #1,\,#2 \right\rangle}
\DeclareMathOperator{\rang}{rang}
\DeclareMathOperator{\diag}{diag}
\DeclareMathOperator{\tr}{tr}
\DeclareMathOperator{\Ker}{ker}
\DeclareMathOperator{\Kep}{Im}
\DeclareMathOperator{\Span}{span}

\title{\LARGE\bfseries Numerikus módszerek\\[6pt]
  \large 12--13.\ hét -- Numerikus integrálás;\
  Szinguláris felbontás és általánosított inverz\\[4pt]
  \normalsize Előadásjegyzet}
\date{2025/26.\ tanév, tavaszi félév}

\begin{document}
\maketitle
\vspace{-1em}
\hrule
\vspace{0.5em}
\tableofcontents
\newpage

% ==============================================================
\part{Numerikus integrálás (12.\ hét)}
% ==============================================================

\section{Az alapfeladat és motiváció}

\subsection*{A feladat}

\textbf{Feladat:} Közelítsük a
\[
  I(f) := \int_a^b f(x)\,dx
\]
Riemann-integrált numerikusan, vagyis véges sok $f$-értékfelhasználásával.

\subsection*{Miért nem a Riemann-definíció?}

A Riemann-integrál definíciója alsó és felső közelítő összegeket (vagy Riemann-összeget) használ.
Ha $[a,b]$-t $m$ egyenlő részre osztjuk, egy tipikus Riemann-összeg
\[
  R_m(f) = \sum_{k=0}^{m-1} f(\xi_k)\,\frac{b-a}{m},
\]
amelyre $R_m(f)\to I(f)$, ha $m\to\infty$. Azonban ez a konvergencia lassú:
az $f$-re tett sima feltételek (például $f\in C^1$) esetén is csak $O(m^{-1})$
sebességű. Tehát kielégítő pontossághoz sok ($m\gg 1$) tagra lenne szükség,
ami számítógépen sok függvényérték-kiértékelést jelent. \textbf{Nem gazdaságos.}

\subsection*{Alkalmazási területek}

A numerikus integrálás szükséges mindazokban az esetekben, amelyekben az analitikus
kiszámítás nem elvégezhető vagy nem kívánatos:
\begin{itemize}[noitemsep,topsep=2pt]
  \item A primitív függvény \emph{nem állítható elő zárt alakban} (például $e^{-x^2}$,
        $\sin(x^2)$, $\cos(x^2)$).
  \item Az analitikus integrálás elvben elvégezhető, de \emph{rendkívül bonyolult}
        lenne.
  \item \emph{Terület, térfogat, ívhossz} kiszámítása, ahol az integrandus mérési
        adatokból adódik.
  \item Differenciálegyenletek numerikus módszereinek konstrukciójánál (például Runge--Kutta
        módszerek, kollokatív módszerek).
\end{itemize}

\begin{mdframed}[backgroundcolor=exbg,linewidth=0pt,innertopmargin=6pt,innerbottommargin=6pt]
\textbf{Motiváló példák:} Számítsuk ki a következő integrálokat!
\[
  \int_0^1 e^{-x^2}\,dx = ?,\qquad \int_0^\pi \cos(x^2)\,dx = ?
\]
Az első a \emph{Gauss-féle hibaintegrálfüggvényhez} kapcsolódik: a primitív függvény nem
elemi. A másodikra szintén nem létezik zárt alakú formula. Mindkettő kiszámítható
numerikus integrálással tetszőleges pontossággal.
\end{mdframed}

\subsection*{Az alaptézis: Lagrange-interpoláció + integráció}

\textbf{Ötlet:} Helyettesítsük $f$-et az $x_0 < x_1 < \cdots < x_n$ alappontokon
épített Lagrange-interpolációs polinomjával $L_n(x) = \sum_{k=0}^n f(x_k)\,\ell_k(x)$,
majd integráljuk az közelített függvényt:
\[
  \int_a^b f(x)\,dx \approx \int_a^b L_n(x)\,dx
  = \int_a^b \sum_{k=0}^n f(x_k)\,\ell_k(x)\,dx
  = \sum_{k=0}^n f(x_k) \underbrace{\int_a^b \ell_k(x)\,dx}_{=:\,A_k}.
\]
Az $A_k := \int_a^b \ell_k(x)\,dx$ \emph{súlyok} csak az alappontoktól (és $[a,b]$-től)
függnek, az $f$ függvénytől \emph{nem}. Ezért elég az alappontokat egyszer rögzíteni,
és utána az $A_k$ értékeket előre kiszámítani. A közelítés:
\[
  \boxed{\int_a^b f(x)\,dx \approx \sum_{k=0}^n A_k\,f(x_k).}
\]

\textbf{Megjegyzés:} az $A_k$ súlyok kizárólag az alappontoktól függnek, $f$-től nem.
Ebből következik, hogy ha $f$ is polinom (és fokszáma legfeljebb $n$), a közelítés
\emph{pontos}: $L_n = f$, tehát $\int_a^b L_n = \int_a^b f$.

% ==============================================================
\section{Kvadratúra formulák általánosan}
\label{sec:kvadratura}

\begin{mdframed}[backgroundcolor=defbg,linewidth=0pt,innertopmargin=8pt,innerbottommargin=8pt]
\begin{definicio}[Kvadratúra formula]
Legyenek $x_0, x_1, \ldots, x_n \in [a,b]$ \emph{alappontok} és $A_0, A_1, \ldots, A_n \in \R$
\emph{súlyok}. Az
\[
  Q_n(f) := \sum_{k=0}^n A_k\,f(x_k)
\]
formulát \emph{kvadratúra formulának} nevezzük.

A kvadratúra formula \emph{interpolációs típusú}, ha $A_k = \int_a^b \ell_k(x)\,dx$
$(k = 0, 1, \ldots, n)$, ahol $\ell_k$ a $k$-adik Lagrange-alappolinom.
\end{definicio}
\end{mdframed}

\begin{mdframed}[backgroundcolor=thmbg,linewidth=0pt,innertopmargin=8pt,innerbottommargin=8pt]
\begin{tetel}[Pontossági tétel]
\label{tetel:pontossagi}
Egy $Q_n$ kvadratúra formula interpolációs típusú $\Longleftrightarrow$ minden
$f \in P_n$ polinomra pontos, azaz
\[
  \int_a^b f(x)\,dx = \sum_{k=0}^n A_k\,f(x_k) \qquad \forall\,f \in P_n.
\]
\end{tetel}
\end{mdframed}

\begin{proof}
\textbf{($\Rightarrow$) Interpolációs típusú $\Rightarrow$ $P_n$-re pontos.}
Ha $f \in P_n$, akkor $L_n(x) = f(x)$ (az interpoláció egyértelműen megadja $f$-et), tehát
\[
  Q_n(f) = \sum_{k=0}^n A_k f(x_k) = \int_a^b \sum_{k=0}^n f(x_k)\,\ell_k(x)\,dx
  = \int_a^b L_n(x)\,dx = \int_a^b f(x)\,dx.
\]

\textbf{($\Leftarrow$) $P_n$-re pontos $\Rightarrow$ interpolációs típusú.}
Alkalmazzuk a feltételt az $f = \ell_j$ Lagrange-alappolinomra ($j = 0, 1, \ldots, n$).
Mivel $\ell_j \in P_n$, a feltétel szerint
\[
  \int_a^b \ell_j(x)\,dx = \sum_{k=0}^n A_k\,\ell_j(x_k) = \sum_{k=0}^n A_k\,\delta_{jk} = A_j.
\]
Tehát $A_j = \int_a^b \ell_j(x)\,dx$, ami pontosan az interpolációs típus definíciója. \qed
\end{proof}

\begin{megjegyzes}
A $Q_n(f) = \sum_{k=0}^n A_k f(x_k)$ formulában összesen $2(n+1)$ szabad paramétert
találunk: az $A_0, \ldots, A_n$ súlyokat és az $x_0, \ldots, x_n$ alappontokat.
Az interpolációs típus megkötése rögzíti az $A_k$-kat (az $x_k$-k függvényeként), így
marad $n+1$ szabad paraméter (az alappontok). Ezek optimális megválasztásával
különféle formulatípusok adódnak:

\begin{enumerate}[noitemsep,topsep=2pt]
  \item \textbf{Newton--Cotes típus:} az alappontok egyenletesen osztják fel $[a,b]$-t.
  \item \textbf{Csebisev típus:} az összes súly azonos: $A_k \equiv A$.
  \item \textbf{Gauss típus:} az alappontokat is optimálisan választják, maximalizálva a
        pontosságot; $2(n+1)$ fokszámig pontos (ez optimális).
\end{enumerate}
A továbbiakban a Newton--Cotes típussal foglalkozunk.
\end{megjegyzes}

% ==============================================================
\section{Newton--Cotes formulák}
\label{sec:nc}

\subsection*{Zárt és nyílt formulák}

A Newton--Cotes formulákban az alappontok egyenletesen helyezkednek el $[a,b]$-ben.

\begin{itemize}[noitemsep,topsep=2pt]
  \item \textbf{Zárt formulák} $Z(n)$: $a$ és $b$ is alappont.
    \[
      h = \frac{b-a}{n}, \quad x_k = a + kh, \quad k = 0, 1, \ldots, n.
    \]
  \item \textbf{Nyílt formulák} $Ny(n)$: $a$ és $b$ \emph{nem} alappont.
    \[
      h = \frac{b-a}{n+2}, \quad x_k = a + (k+1)h, \quad k = 0, 1, \ldots, n.
    \]
\end{itemize}

\subsection*{Az együtthatók meghatározása}

A pontossági tétel alapján az $A_k$ értékeket a következő lineáris egyenletrendszer adja meg:
\[
  \int_a^b x^j\,dx = \sum_{k=0}^n A_k\,x_k^j, \qquad j = 0, 1, \ldots, n.
\]
Explicit alakban:
\[
  \frac{b^{j+1} - a^{j+1}}{j+1} = A_0 x_0^j + A_1 x_1^j + \cdots + A_n x_n^j,
  \qquad j = 0, 1, \ldots, n.
\]
A rendszer mátrixa
\[
  V^T =
  \begin{pmatrix}
    1 & 1 & \cdots & 1 \\
    x_0 & x_1 & \cdots & x_n \\
    \vdots & \vdots & & \vdots \\
    x_0^n & x_1^n & \cdots & x_n^n
  \end{pmatrix},
\]
azaz a Vandermonde-mátrix transzponáltja; ez nem szinguláris, ha az alappontok különbözők.

\subsection*{Az érintő formula (Ny(0))}

A nyílt $n=0$ esetben egyetlen alappont van, a középpont:
$x_0 = a + h = a + \tfrac{b-a}{2} = \tfrac{a+b}{2}$.
A LER egyetlen egyenlete: $A_0 \cdot 1 = b - a$, tehát $A_0 = b-a$.

\begin{mdframed}[backgroundcolor=thmbg,linewidth=0pt,innertopmargin=8pt,innerbottommargin=8pt]
\textbf{Érintő formula} (középponti szabály, $Ny(0)$):
\[
  \int_a^b f(x)\,dx \approx (b-a)\cdot f\!\left(\frac{a+b}{2}\right) =: E(f).
\]
\end{mdframed}

\emph{Geometriai értelmezés:} a görbe alatti területet azzal a téglalappal közelítjük,
amelynek magassága $f$ értéke a középpontban. Ez nem tényleges érintőegyenesre, hanem
a középpontban vett vízszintes magasságra támaszkodik. Az „érintő formula” elnevezés
történeti okokból maradt használatban; ma gyakran középponti szabálynak nevezzük.

\subsection*{A trapéz formula (Z(1))}

A zárt $n=1$ esetben $x_0 = a$, $x_1 = b$.
A LER:
\[
  A_0 + A_1 = b - a, \qquad A_0 \cdot a + A_1 \cdot b = \tfrac{b^2-a^2}{2} = \tfrac{(b-a)(b+a)}{2}.
\]
Az első egyenletből $A_0 = b-a-A_1$; visszahelyettesítve: $(b-a-A_1)a + A_1 b = \tfrac{(b-a)(a+b)}{2}$,
ami $A_1(b-a) = \tfrac{(b-a)(a+b)}{2} - a(b-a) = \tfrac{(b-a)^2}{2}$, tehát $A_1 = \tfrac{b-a}{2}$.
Szimmetria alapján $A_0 = \tfrac{b-a}{2}$.

\begin{mdframed}[backgroundcolor=thmbg,linewidth=0pt,innertopmargin=8pt,innerbottommargin=8pt]
\textbf{Trapéz formula} ($Z(1)$):
\[
  \int_a^b f(x)\,dx \approx \frac{b-a}{2}\cdot\bigl(f(a) + f(b)\bigr) =: T(f).
\]
\end{mdframed}

\emph{Geometriai értelmezés:} a görbe alatti területet a $(a,0)$, $(a,f(a))$, $(b,f(b))$,
$(b,0)$ csúcsú trapéz területével helyettesítjük.

\subsection*{A Simpson formula (Z(2))}

A zárt $n=2$ esetben $x_0=a$, $x_1=\tfrac{a+b}{2}$, $x_2=b$, $h=\tfrac{b-a}{2}$.
A LER:
\begin{align*}
  A_0 + A_1 + A_2 &= b - a, \\
  A_0 a + A_1 \tfrac{a+b}{2} + A_2 b &= \tfrac{b^2-a^2}{2}, \\
  A_0 a^2 + A_1 \left(\tfrac{a+b}{2}\right)^2 + A_2 b^2 &= \tfrac{b^3-a^3}{3}.
\end{align*}
Szimmetria alapján $A_0 = A_2$. Az első egyenletből $2A_0 + A_1 = b-a$.
A megoldás: $A_0 = A_2 = \tfrac{b-a}{6}$, $A_1 = \tfrac{4(b-a)}{6} = \tfrac{2(b-a)}{3}$.

\begin{mdframed}[backgroundcolor=thmbg,linewidth=0pt,innertopmargin=8pt,innerbottommargin=8pt]
\textbf{Simpson formula} ($Z(2)$):
\[
  \int_a^b f(x)\,dx \approx \frac{b-a}{6}\left(f(a) + 4f\!\left(\frac{a+b}{2}\right)
  + f(b)\right) =: S(f).
\]
\end{mdframed}

\emph{Megjegyzés:} A Simpson-formula a konstrukció alapján $P_2$-re pontos.
Ennél azonban \emph{erősebb}: $P_3$-ra is pontos. Ennek oka a középpontra vett
szimmetria. Ha $c=\tfrac{a+b}{2}$, akkor minden harmadfokú polinom felírható
\[
  p(x)=\alpha_0+\alpha_1(x-c)+\alpha_2(x-c)^2+\alpha_3(x-c)^3
\]
alakban. Az $(x-c)$ és $(x-c)^3$ tagok integrálja a szimmetrikus $[c-h,c+h]$
intervallumon nulla, és a Simpson-formula is nullát ad rájuk. A konstans és másodfokú
tagokra a formula már az interpolációs konstrukció miatt pontos. Így a Simpson-formula
minden $p\in P_3$ polinomra pontos. Például $p(x)=x^3$ esetén közvetlenül is ellenőrizhető, hogy
\[
S(x^3)=\frac{b-a}{6}\left(a^3+4\left(\frac{a+b}{2}\right)^3+b^3\right)=\frac{b^4-a^4}{4}=
\int_a^b x^3\,dx. \quad \checkmark
\]

Az ábrák geometriai szemléltetést nyújtanak:

\begin{figure}[h!]
\centering
\includegraphics[width=0.92\textwidth]{quad_formulak.pdf}
\caption{Az érintő, a trapéz és a Simpson formula szemléltetése ($f(x)=\sin x$, $[0{,}4;\;2{,}8]$).
  Az aranyszínű terület a közelítő terület, a kék görbe az eredeti függvény.}
\label{fig:quad}
\end{figure}

% ==============================================================
\section{Hibaformulák}
\label{sec:hiba}

A hibaformulák megadják, mekkora eltérés keletkezik az egzakt és a közelítő integrálérték között.
Az alapeszköz az integrálszámítás középértéktétele.

\begin{mdframed}[backgroundcolor=thmbg,linewidth=0pt,innertopmargin=8pt,innerbottommargin=8pt]
\begin{tetel}[Az integrálszámítás középértéktétele]
\label{tetel:kvmvt}
Ha $f \in C[a,b]$ és $g \in C[a,b]$, továbbá $g(x) \ge 0$ minden $x \in [a,b]$-re, akkor
létezik olyan $\xi \in [a,b]$, hogy
\[
  \int_a^b f(x)\,g(x)\,dx = f(\xi)\int_a^b g(x)\,dx.
\]
\end{tetel}
\end{mdframed}

\begin{proof}
Legyen $m = \min_{[a,b]} f$ és $M = \max_{[a,b]} f$. Ekkor
$m\,g(x) \le f(x)\,g(x) \le M\,g(x)$, ezért integrálással
\[
  m\int_a^b g(x)\,dx \le \int_a^b f(x)g(x)\,dx \le M\int_a^b g(x)\,dx.
\]
Ha $\int_a^b g(x)\,dx = 0$, akkor $g \ge 0$ miatt $g \equiv 0$, így az állítás triviális.
Ellenkező esetben
\[
  m \le \frac{\int_a^b f(x)g(x)\,dx}{\int_a^b g(x)\,dx} \le M.
\]
A folytonosság miatt $f$ felveszi az $m$ és $M$ közötti összes értéket, tehát létezik
$\xi \in [a,b]$, amelyre
\[
  f(\xi)=\frac{\int_a^b f(x)g(x)\,dx}{\int_a^b g(x)\,dx}.
\]
Ebből következik az állítás. \qed
\end{proof}

\subsection*{Az érintő formula hibája}

\begin{mdframed}[backgroundcolor=thmbg,linewidth=0pt,innertopmargin=8pt,innerbottommargin=8pt]
\begin{tetel}
\label{tetel:erintohiba}
Ha $f \in C^2[a,b]$, akkor $\exists\,\eta \in [a,b]$:
\[
  \int_a^b f(x)\,dx - E(f) = \frac{(b-a)^3}{24}\,f''(\eta).
\]
\end{tetel}
\end{mdframed}

\begin{proof}
Legyen $c = \tfrac{a+b}{2}$ és $h = \tfrac{b-a}{2}$. Ekkor
\[
  \int_a^b f(x)\,dx - E(f)
  = \int_0^h \bigl(f(c+t)+f(c-t)-2f(c)\bigr)\,dt.
\]
Taylor integrálmaradékos alakját használva
\[
  f(c+t)-f(c)-t f'(c)=\int_0^t (t-s)f''(c+s)\,ds,
\]
és
\[
  f(c-t)-f(c)+t f'(c)=\int_0^t (t-s)f''(c-s)\,ds.
\]
A két egyenletet összeadva, majd $t$ szerint $0$-tól $h$-ig integrálva:
\[
  \int_a^b f(x)\,dx - E(f)
  = \int_0^h \frac{(h-s)^2}{2}\bigl(f''(c+s)+f''(c-s)\bigr)\,ds.
\]
A súlyfüggvény nemnegatív, $f''$ folytonos, ezért az integrálszámítás középértéktétele
szerint létezik $\eta\in[a,b]$, hogy
\[
  \int_a^b f(x)\,dx - E(f)
  = f''(\eta)\int_0^h (h-s)^2\,ds
  = \frac{h^3}{3}f''(\eta)
  = \frac{(b-a)^3}{24}f''(\eta). \qed
\]
\end{proof}

\subsection*{A trapéz formula hibája}

\begin{mdframed}[backgroundcolor=thmbg,linewidth=0pt,innertopmargin=8pt,innerbottommargin=8pt]
\begin{tetel}
\label{tetel:trapezhiba}
Ha $f \in C^2[a,b]$, akkor $\exists\,\eta \in [a,b]$:
\[
  \int_a^b f(x)\,dx - T(f) = -\frac{(b-a)^3}{12}\,f''(\eta).
\]
\end{tetel}
\end{mdframed}

\begin{proof}
Legyen $L_1$ az $f$ függvény $a$ és $b$ pontokban vett lineáris interpolációs polinomja.
Ekkor $T(f)=\int_a^b L_1(x)\,dx$. A lineáris interpoláció maradéktagja alapján minden
$x\in(a,b)$-re létezik $\xi_x\in(a,b)$, amelyre
\[
  f(x)-L_1(x)=\frac{f''(\xi_x)}{2}(x-a)(x-b).
\]
Mivel $(x-a)(x-b)\le 0$ az $[a,b]$ intervallumon, a szokásos súlyozott középértékes
hibaalakból kapjuk, hogy létezik $\eta\in[a,b]$, amelyre
\[
  \int_a^b f(x)\,dx-T(f)
  = \frac{f''(\eta)}{2}\int_a^b (x-a)(x-b)\,dx.
\]
Az utolsó integrál
\[
  \int_a^b (x-a)(x-b)\,dx
  = -\int_a^b (x-a)(b-x)\,dx
  = -\frac{(b-a)^3}{6}.
\]
Ezért
\[
  \int_a^b f(x)\,dx - T(f) = -\frac{(b-a)^3}{12}f''(\eta). \qed
\]
\end{proof}

\begin{megjegyzes}
Az érintő és a trapéz formula hibáját összehasonlítva:
\[
  \int_a^b f - E(f) \approx \frac{(b-a)^3}{24}f''(\eta),\quad
  \int_a^b f - T(f) \approx -\frac{(b-a)^3}{12}f''(\eta).
\]
Az érintő formula hibája körülbelül \emph{felakkora, és ellentétes előjelű}, mint a trapézé:
$E$-hiba $\approx -\tfrac{1}{2}\cdot T$-hiba. Ez az összefüggés a Richardson-extrapoláció
alapja: $\tfrac{2E(f)+T(f)}{3}$ a Simpson-formula!

Valóban: $\tfrac{1}{3}T(f) + \tfrac{2}{3}E(f) = \tfrac{b-a}{6}(f(a)+f(b)) + \tfrac{2(b-a)}{3}f(c) = S(f)$.
\end{megjegyzes}

\subsection*{A Simpson formula hibája}

\begin{mdframed}[backgroundcolor=thmbg,linewidth=0pt,innertopmargin=8pt,innerbottommargin=8pt]
\begin{tetel}
\label{tetel:simpsonhiba}
Ha $f \in C^4[a,b]$, akkor $\exists\,\eta \in [a,b]$:
\[
  \int_a^b f(x)\,dx - S(f) = -\frac{(b-a)^5}{2880}\,f^{(4)}(\eta).
\]
\end{tetel}
\end{mdframed}

\begin{proof}
A Simpson-formula minden legfeljebb harmadfokú polinomra pontos, ezért a hibafunkcionál
\[
  L(f):=\int_a^b f(x)\,dx-S(f)
\]
eltűnik $P_3$-on. A Peano-féle hibatag-tétel szerint, ha $f\in C^4[a,b]$, akkor
\[
  L(f)=\int_a^b K(t)f^{(4)}(t)\,dt
\]
egy alkalmas Peano-kernellel. A Simpson-formula szimmetriája miatt ez a kernel előjeltartó,
így a súlyozott középértéktétel alkalmazható: létezik $\eta\in[a,b]$, amelyre
\[
  L(f)=f^{(4)}(\eta)\,L\!\left(\frac{x^4}{4!}\right).
\]
A konstans meghatározásához elegendő az $f(x)=x^4$ esetet kiszámítani. A skálázás miatt
feltehetjük, hogy $[a,b]=[-h,h]$. Ekkor
\[
  \int_{-h}^{h} x^4\,dx=\frac{2h^5}{5},
  \qquad
  S(x^4)=\frac{h}{3}\bigl(h^4+4\cdot 0+h^4\bigr)=\frac{2h^5}{3}.
\]
Tehát
\[
  L(x^4)=\frac{2h^5}{5}-\frac{2h^5}{3}=-\frac{4h^5}{15},
  \qquad
  L\!\left(\frac{x^4}{4!}\right)=-\frac{h^5}{90}.
\]
Mivel $h=\tfrac{b-a}{2}$,
\[
  \int_a^b f(x)\,dx-S(f)
  =-\frac{h^5}{90}f^{(4)}(\eta)
  =-\frac{(b-a)^5}{2880}f^{(4)}(\eta). \qed
\]
\end{proof}

\begin{mdframed}[backgroundcolor=exbg,linewidth=0pt,innertopmargin=6pt,innerbottommargin=6pt]
\textbf{Összefoglalás: egyszerű formulák hibái}
\[
\renewcommand{\arraystretch}{1.5}
\begin{array}{l|c|c}
\text{Formula} & \text{Hibatag} & \text{Pontossági fok}\\
\hline
E(f)\text{ (érintő)} & +\dfrac{(b-a)^3}{24}\,f''(\eta) & P_1 \\
T(f)\text{ (trapéz)} & -\dfrac{(b-a)^3}{12}\,f''(\eta) & P_1 \\
S(f)\text{ (Simpson)} & -\dfrac{(b-a)^5}{2880}\,f^{(4)}(\eta) & P_3 \\
\end{array}
\]
\end{mdframed}

% ==============================================================
\section{Összetett formulák}
\label{sec:osszetett}

\subsection*{Motiváció: miért szükségesek?}

A hibák $(b-a)^3$-ra illetve $(b-a)^5$-re arányosak. Ha $[a,b]$ nagy, a hiba
is nagy lehet. A megoldás: \emph{osszuk fel $[a,b]$-t sok kis részre}, és minden
részintervallumon alkalmazzuk az egyszerű formulát. Az összeg teljes hibáját
a kis részhibák összege adja.

\subsection*{A trapéz összetett formula}

Osszuk fel $[a,b]$-t $m$ egyenlő részre: $h = \tfrac{b-a}{m}$,
$x_k = a + kh$ $(k = 0, 1, \ldots, m)$.
Alkalmazzuk a trapéz formulát minden $[x_{k-1}, x_k]$ részintervallumon:
\[
  \int_{x_{k-1}}^{x_k} f \approx \frac{h}{2}\bigl(f(x_{k-1}) + f(x_k)\bigr).
\]
Összegezve $k = 1, \ldots, m$-re, és felhasználva, hogy a belső pontok $(x_1, \ldots, x_{m-1})$
kétszer jelennek meg:

\begin{mdframed}[backgroundcolor=thmbg,linewidth=0pt,innertopmargin=8pt,innerbottommargin=8pt]
\textbf{Trapéz összetett formula} (trapéz szabály):
\[
  \int_a^b f(x)\,dx \approx T_m(f) := \frac{b-a}{2m}
  \left(f(a) + 2\sum_{k=1}^{m-1} f(x_k) + f(b)\right).
\]
Az együtthatósorozat: $1,\,2,\,2,\,\ldots,\,2,\,1$.
\end{mdframed}

\begin{mdframed}[backgroundcolor=thmbg,linewidth=0pt,innertopmargin=8pt,innerbottommargin=8pt]
\begin{tetel}[A trapéz összetett formula hibája]
\label{tetel:trapezosszetett}
Ha $f \in C^2[a,b]$, akkor $\exists\,\eta \in [a,b]$:
\[
  \int_a^b f(x)\,dx - T_m(f) = -\frac{(b-a)^3}{12m^2}\,f''(\eta).
\]
\end{tetel}
\end{mdframed}

\begin{proof}
Az $i$-edik részintervallumon ($i = 1, \ldots, m$) a trapéz hiba
$-\tfrac{h^3}{12}f''(\eta_i)$, ahol $h = \tfrac{b-a}{m}$ és
$\eta_i\in[x_{i-1},x_i]$. Összegezve:
\[
  \int_a^b f - T_m(f) = -\frac{h^3}{12}\sum_{i=1}^m f''(\eta_i).
\]
Mivel $f''$ folytonos, a véges sok érték számtani közepe is $f''$ minimuma és maximuma
közé esik; ezért létezik $\eta\in[a,b]$, amelyre
\[
  \frac1m\sum_{i=1}^m f''(\eta_i)=f''(\eta).
\]
Így
\[
  \int_a^b f - T_m(f)
  = -\frac{m h^3}{12} f''(\eta)
  = -\frac{(b-a)^3}{12m^2}f''(\eta). \qed
\]
\end{proof}

\subsection*{A Simpson összetett formula}

A Simpson formula páros $m$-t igényel ($m = 2s$, $s \in \N$). Ekkor $[a,b]$-t $m$ egyenlő
részre osztjuk ($h = \tfrac{b-a}{m}$), majd a $[x_{2k-2}, x_{2k}]$ ($k = 1, \ldots, m/2$)
részintervallumokon -- amelyek szélessége $2h$ -- alkalmazzuk a Simpson formulát.

\begin{mdframed}[backgroundcolor=thmbg,linewidth=0pt,innertopmargin=8pt,innerbottommargin=8pt]
\textbf{Simpson összetett formula} (Simpson szabály):
\[
  S_m(f) := \frac{b-a}{3m}
  \left(f(a) + 4\sum_{k=1}^{m/2} f(x_{2k-1}) + 2\sum_{k=1}^{m/2-1} f(x_{2k}) + f(b)\right).
\]
Az együtthatósorozat: $1,\,4,\,2,\,4,\,2,\,\ldots,\,4,\,2,\,4,\,1$.
\end{mdframed}

\begin{mdframed}[backgroundcolor=thmbg,linewidth=0pt,innertopmargin=8pt,innerbottommargin=8pt]
\begin{tetel}[A Simpson összetett formula hibája]
\label{tetel:simpsonosszetett}
Ha $f \in C^4[a,b]$, akkor $\exists\,\eta \in [a,b]$:
\[
  \int_a^b f(x)\,dx - S_m(f) = -\frac{(b-a)^5}{180\,m^4}\,f^{(4)}(\eta).
\]
\end{tetel}
\end{mdframed}

\begin{proof}
Az egyes Simpson-panelek szélessége $2h$, ahol $h = \tfrac{b-a}{m}$.
Az $i$-edik panelen a Simpson-hiba
$-\tfrac{(2h)^5}{2880}f^{(4)}(\eta_i)$ valamely panelbeli $\eta_i$-re.
Összesen $m/2$ panel van, ezért
\[
  \int_a^b f - S_m(f)
  = -\frac{(2h)^5}{2880}\sum_{i=1}^{m/2} f^{(4)}(\eta_i).
\]
A $f^{(4)}$ folytonossága miatt a véges sok érték átlaga felvehető egy
$\eta\in[a,b]$ pontban, így
\[
  \int_a^b f - S_m(f)
  = -\frac{m}{2}\cdot\frac{32h^5}{2880}f^{(4)}(\eta)
  = -\frac{16m h^5}{2880}f^{(4)}(\eta)
  = -\frac{(b-a)^5}{180m^4}f^{(4)}(\eta). \qed
\]
\end{proof}

\begin{megjegyzes}
\leavevmode
\begin{enumerate}[noitemsep,topsep=2pt]
  \item Ha $f \in C^2[a,b]$, akkor $T_m(f) \to I(f)$, és a konvergencia $m^2$ rendű.
  \item Ha $f \in C^4[a,b]$, akkor $S_m(f) \to I(f)$, és a konvergencia $m^4$ rendű.
  \item A Simpson összegzett formula tehát lényegesen gyorsabban konvergál, mint a trapéz.
  \item Az érintő formulából is képezhető összetett formula; hibája szintén $O(m^{-2})$,
        de az előjele ellentétes a trapézéval.
\end{enumerate}
\end{megjegyzes}

\begin{figure}[h!]
\centering
\includegraphics[width=0.95\textwidth]{osszetett_trapez.pdf}
\caption{Az összetett trapéz formula $m = 1, 2, 4, 8$ résszel ($f(x)=\sin x$, $[0,\pi]$).
  Az aranyszínű terület a közelítő, a kék görbe az egzakt függvény.}
\label{fig:osszetett}
\end{figure}

\begin{mdframed}[backgroundcolor=exbg,linewidth=0pt,innertopmargin=8pt,innerbottommargin=8pt]
\textbf{Számpélda.} Közelítsük az $\int_0^1 e^{-x^2}\,dx$ integrált $m = 4$ résszel,
a trapéz és a Simpson összetett formulájával!

Az osztópontok: $h = 0{,}25$; $x_k = 0;\;0{,}25;\;0{,}5;\;0{,}75;\;1{,}0$.
Függvényértékek: $f_k = e^{-x_k^2}$:
\[
  f_0 = 1{,}00000,\quad f_1 = 0{,}93941,\quad f_2 = 0{,}77880,\quad
  f_3 = 0{,}57007,\quad f_4 = 0{,}36788.
\]

\textbf{Trapéz $(m=4)$:}
$T_4 = \tfrac{0{,}25}{2}(f_0 + 2f_1 + 2f_2 + 2f_3 + f_4)
= \tfrac{1}{8}(1 + 2\cdot 0{,}93941 + 2\cdot 0{,}77880 + 2\cdot 0{,}57007 + 0{,}36788)
= \tfrac{1}{8}\cdot 5{,}94444 \approx 0{,}74306.$

\textbf{Simpson $(m=4)$:}
$S_4 = \tfrac{0{,}25}{3}(f_0 + 4f_1 + 2f_2 + 4f_3 + f_4)
= \tfrac{1}{12}(1 + 4\cdot 0{,}93941 + 2\cdot 0{,}77880 + 4\cdot 0{,}57007 + 0{,}36788)
\approx 0{,}74683.$

A pontos érték (12 jegy): $0{,}74682413\ldots$
Hiba: $T_4$-nél $\approx 3{,}8\cdot 10^{-3}$, $S_4$-nél $\approx 7\cdot 10^{-7}$. A Simpson
formula négy részre osztva is több nagyságrenddel pontosabb!
\end{mdframed}

% ==============================================================
\part{Szinguláris felbontás és általánosított inverz (13.\ hét)}
% ==============================================================

\section{Szinguláris felbontás (SVD)}
\label{sec:svd}

\subsection*{Motiváció}

Hermite-féle (valós esetben szimmetrikus) $A \in \C^{n\times n}$ mátrix esetén az $A = UDU^*$ spektrális felbontás létezik,
ahol $U$ unitér és $D = \diag(\lambda_1, \ldots, \lambda_n)$. Ez azt fejezi ki: $A$ az
az $u_1, \ldots, u_n$ ortonormált bázis vektorait sajátvektorokként kezeli, azaz $Au_i=\lambda_i u_i$.

Nem szimmetrikus esetben ez általában nem lehetséges. Azonban azt elérhetjük, hogy $A$ az
$\R^n$ egy ortonormált bázisát ($v_1, \ldots, v_r$) $\R^m$ egy másik ortonormált rendszerébe
($u_1, \ldots, u_r$) vigye, $\sigma_i > 0$ skalárok segítségével:
$A v_i = \sigma_i u_i$.

\begin{mdframed}[backgroundcolor=thmbg,linewidth=0pt,innertopmargin=8pt,innerbottommargin=8pt]
\begin{tetel}[Szinguláris felbontás]
\label{tetel:svd}
Tetszőleges $A \in \C^{m\times n}$ mátrixhoz $\exists\, U \in \C^{m\times m}$ és
$V \in \C^{n\times n}$ unitér mátrix, valamint $D \in \C^{m\times n}$ mátrix, amelyre
\[
  A = UDV^*,
\]
ahol
\[
  D = \begin{pmatrix} \Sigma & 0 \\ 0 & 0 \end{pmatrix}, \qquad
  \Sigma = \diag(\sigma_1, \sigma_2, \ldots, \sigma_r), \quad
  \sigma_1 \ge \sigma_2 \ge \cdots \ge \sigma_r > 0,
\]
és $r = \rang(A)$.
\end{tetel}
\end{mdframed}

\begin{proof}
\textbf{Az $A^*A$ mátrix vizsgálata.}
$A^*A \in \C^{n\times n}$ Hermite-mátrix ($(A^*A)^* = A^*(A^*)^* = A^*A$) és pozitív szemidefinit:
$x^*(A^*A)x = (Ax)^*(Ax) = \norm{Ax}^2 \ge 0$. Ezért $A^*A$ unitéren diagonalizálható:
$A^*A = V \Lambda V^*$, ahol $V$ unitér és $\Lambda = \diag(\lambda_1, \ldots, \lambda_n)$,
$\lambda_1 \ge \cdots \ge \lambda_n \ge 0$.

A nullánál nagyobb sajátértékek száma $r = \rang(A^*A) = \rang(A)$.

\textbf{A szinguláris értékek.}
Legyen $\sigma_i = \sqrt{\lambda_i} > 0$ $(i = 1, \ldots, r)$, és legyenek $v_1, \ldots, v_n$
$V$ unitér mátrix oszlopai ($A^*A$ ortonormált sajátvektorai).

\textbf{A bal szinguláris vektorok meghatározása.}
Definiáljuk: $u_i := \dfrac{Av_i}{\sigma_i}$ $(i = 1, \ldots, r)$. Ellenőrizzük az
ortogonalitást: $i \ne j$ esetén
\[
  \innerprod{u_i}{u_j} = \frac{1}{\sigma_i\sigma_j}\innerprod{Av_i}{Av_j}
  = \frac{1}{\sigma_i\sigma_j} v_i^*(A^*A)v_j
  = \frac{\lambda_j}{\sigma_i\sigma_j}\innerprod{v_i}{v_j} = 0.
\]
Ugyanígy $\norm{u_i}^2 = \tfrac{1}{\sigma_i^2}\norm{Av_i}^2 = \tfrac{\lambda_i}{\sigma_i^2} = 1$.
Tehát $u_1, \ldots, u_r$ ortonormált rendszer. Egészítsük ki $\C^m$ ortonormált bázisává:
$u_1, \ldots, u_m$.

\textbf{Az $A = UDV^*$ egyenlet igazolása.}
$AV = A[v_1 \mid \cdots \mid v_n] = [\sigma_1 u_1 \mid \cdots \mid \sigma_r u_r \mid 0 \mid \cdots \mid 0]
= U\,D,$
ezért $A = UD\,V^*$ (jobb oldalon $V^*$-gal szorzunk, mivel $V$ unitér). \qed
\end{proof}

\begin{mdframed}[backgroundcolor=defbg,linewidth=0pt,innertopmargin=8pt,innerbottommargin=8pt]
\begin{definicio}[Szinguláris értékek]
Az $A \in \C^{m\times n}$ mátrix szinguláris felbontásában a $\sigma_1, \ldots, \sigma_r > 0$
értékeket az $A$ \emph{szinguláris értékeinek} nevezzük. Felírhatók mint
\[
  \sigma_i = \sqrt{\lambda_i(A^*A)} > 0, \qquad i = 1, \ldots, r.
\]
A $v_i$ vektorokat \emph{jobb szinguláris vektoroknak}, az $u_i$ vektorokat
\emph{bal szinguláris vektoroknak} nevezzük.
\end{definicio}
\end{mdframed}

\begin{megjegyzes}
\leavevmode
\begin{enumerate}[noitemsep,topsep=2pt]
  \item $A^*A \in \C^{n\times n}$ és $AA^* \in \C^{m\times m}$ pozitív sajátértékei
        \emph{azonosak}: ha $A^*Av = \lambda v$, akkor $AA^*(Av) = \lambda(Av)$.
  \item $A^*A$ sajátvektorai a $V$ unitér mátrix oszlopai (jobb szinguláris vektorok).
  \item $AA^*$ sajátvektorai az $U$ unitér mátrix oszlopai (bal szinguláris vektorok).
  \item $m = n$ esetén: $\sigma_1 = \norm{A}_2$ (indukált 2-norma) és
        $\norm{A}_F^2 = \sum_{i=1}^r \sigma_i^2$ (Frobenius-norma).
\end{enumerate}
\end{megjegyzes}

\begin{mdframed}[backgroundcolor=thmbg,linewidth=0pt,innertopmargin=8pt,innerbottommargin=8pt]
\begin{tetel}[Lineáris transzformáció jellemzése]
\label{tetel:lintrafo}
Ha $\rang(A) = r$, akkor
\[
  \Ker(A) = \Span\{v_{r+1}, \ldots, v_n\}, \qquad
  \Kep(A) = \Span\{u_1, \ldots, u_r\}.
\]
\end{tetel}
\end{mdframed}

\begin{proof}
\textbf{Ker:} Ha $v \in \Ker(A)$, akkor $A^*Av = 0$, azaz $v$ a nulla sajátértékhez tartozó
sajátvektor: $v \in \Span\{v_{r+1},\ldots,v_n\}$. Megfordítva: $Av_j = 0$ minden $j > r$-re,
mert $\norm{Av_j}^2 = \lambda_j = 0$.

\textbf{Im:} $\Kep(A) = \Span\{Av_1,\ldots,Av_r\} = \Span\{\sigma_1 u_1,\ldots,\sigma_r u_r\}
= \Span\{u_1,\ldots,u_r\}$. \qed
\end{proof}

\subsection*{Alkalmazási területek}

\begin{enumerate}[noitemsep,topsep=2pt]
  \item \textbf{Többváltozós statisztika:} a szóródási mátrix vizsgálata, varianciák tömörítése
        (főkomponens-analízis, PCA).
  \item \textbf{Dokumentumkeresés és tematizálás:} nagy adatbázisokban latens szemantikus
        analízis (LSA), adatbányászat.
  \item \textbf{Jelfeldolgozás:} képtömörítés, zajszűrés, alacsony rangú közelítések.
  \item \textbf{Geofizika:} mérési adatok inverziója, tomográfia.
\end{enumerate}

% ==============================================================
\section{Az SVD alkalmazása: Eckart--Young-tétel és képtömörítés}
\label{sec:eckartyoung}

\subsection*{Alacsony rangú közelítés}

A szinguláris felbontásból $A = \sum_{i=1}^r \sigma_i u_i v_i^*$ következik
(az SVD kifejtett alakja). Legyen $A_k = \sum_{i=1}^k \sigma_i u_i v_i^*$ az első $k$ tag.
Ekkor $\rang(A_k) = k$.

\begin{mdframed}[backgroundcolor=thmbg,linewidth=0pt,innertopmargin=8pt,innerbottommargin=8pt]
\begin{tetel}[Eckart--Young]
\label{tetel:ey}
Legyen $\rang(A) = r$ és $0\le k<r$; továbbá
$A_k=\sum_{i=1}^k \sigma_i u_i v_i^*$ ($A_0=0$). Ekkor $A_k$ az $A$ legjobb
rangfeljebb-$k$ közelítése a 2-normában \emph{és} a Frobenius-normában egyaránt:
\[
  \norm{A - A_k}_F = \sqrt{\sum_{i=k+1}^r \sigma_i^2}, \qquad
  \norm{A - A_k}_2 = \sigma_{k+1}.
\]
Vagyis minden $\rang(B) \le k$ mátrixra
\[
  \norm{A-B}_2\ge \sigma_{k+1},\qquad
  \norm{A-B}_F\ge \sqrt{\sum_{i=k+1}^r\sigma_i^2}.
\]
\end{tetel}
\end{mdframed}

\begin{proof}
Először vegyük észre, hogy
$A - A_k = \sum_{i=k+1}^r \sigma_i u_i v_i^*$, ezért
\[
  \norm{A-A_k}_2=\sigma_{k+1},\qquad
  \norm{A-A_k}_F=\sqrt{\sum_{i=k+1}^r\sigma_i^2}.
\]
A 2-normás optimalitáshoz legyen $B$ tetszőleges, $\rang(B)\le k$ mátrix, és
$W = \Ker(B)$. Ekkor $\dim(W) \ge n-k$. A
$V_{k+1} = \Span\{v_1,\ldots,v_{k+1}\}$ altér dimenziója $k+1$, ezért
$W\cap V_{k+1}$ tartalmaz nemnulla vektort. Válasszunk benne egységvektort:
$w = \sum_{i=1}^{k+1} c_i v_i$, $\sum |c_i|^2 = 1$. Mivel $Bw=0$,
\[
  \norm{A - B}_2 \ge \norm{(A-B)w}_2 = \norm{Aw}_2
  = \left\|\sum_{i=1}^{k+1} c_i \sigma_i u_i\right\|_2
  = \sqrt{\sum_{i=1}^{k+1} |c_i|^2 \sigma_i^2} \ge \sigma_{k+1}.
\]
A Frobenius-normás optimalitás a Mirsky-tétel speciális esete: a Frobenius-norma
unitér invarianciája miatt a legjobb rangfeljebb-$k$ közelítés hibája pontosan a
levágott szinguláris értékek négyzetösszegének gyöke. Ezt az $A_k$ mátrix eléri,
tehát a megadott alsó korlát éles. \qed
\end{proof}

\subsection*{Alkalmazás: képtömörítés}

Egy szürkeárnyalatos pixelkép reprezentálható mint $A \in \R^{m\times n}$ mátrix, ahol
$a_{ij}$ a $(i,j)$ pixel intenzitása. Az SVD felbontás után:
\begin{itemize}[noitemsep,topsep=2pt]
  \item Az $A_k$ rang-$k$ közelítés \emph{tömörített} képet ad: tároláshoz $k(m+n)$ szám
        kell (az $\sigma_i$, $u_i$, $v_i$ adatok), szemben az eredeti $mn$-nel.
  \item A tömörítési arány akkor jó, ha a szinguláris értékek gyorsan csökkennek
        (néhány nagy dominálja a többit).
  \item A Frobenius-hiba: $\norm{A - A_k}_F = \sqrt{\sum_{i=k+1}^r \sigma_i^2}$,
        amely a pixelenkénti négyzethibát méri.
\end{itemize}

\begin{figure}[h!]
\centering
\includegraphics[width=0.88\textwidth]{svd_tomorites.pdf}
\caption{SVD-alapú alacsony rangú közelítés egy szintetikus képen.
  Bal: az eredeti (teljes rangú) kép. Közép: rang-$r/2$ közelítés. Jobb: rang-10 közelítés.}
\label{fig:svd}
\end{figure}

% ==============================================================
\section{Az általánosított inverz (Moore--Penrose-pszeudoinverz)}
\label{sec:pseudoinv}

\subsection*{Motiváció}

Ha $A \in \C^{m\times n}$ nem négyzetes ($m \ne n$), vagy négyzetes de szinguláris, akkor
az $Ax = b$ LER általában nem oldható meg klasszikus értelemben. Keresünk egy
``legjobb közelítő megoldást'', amelyet az $A^+$ pszeudoinverz segítségével írhatunk fel.

\begin{mdframed}[backgroundcolor=defbg,linewidth=0pt,innertopmargin=8pt,innerbottommargin=8pt]
\begin{definicio}[Moore--Penrose-féle általánosított inverz]
Az $A \in \C^{m\times n}$ mátrix \emph{általánosított (pszeudo-) inverze} az $A^+ \in \C^{n\times m}$
mátrix, ha a következő négy \emph{Moore--Penrose-feltétel} teljesül:
\begin{enumerate}[noitemsep,topsep=2pt]
  \item $AA^+$ önadjungált: $(AA^+)^* = AA^+$,
  \item $A^+A$ önadjungált: $(A^+A)^* = A^+A$,
  \item $AA^+A = A$,
  \item $A^+AA^+ = A^+$.
\end{enumerate}
\end{definicio}
\end{mdframed}

\begin{mdframed}[backgroundcolor=thmbg,linewidth=0pt,innertopmargin=8pt,innerbottommargin=8pt]
\begin{tetel}
Az $A^+$ általánosított inverz \emph{egyértelmű}.
\end{tetel}
\end{mdframed}

\begin{proof}
Tegyük fel, hogy $B$ és $C$ is teljesíti a Moore--Penrose-feltételeket $A$-ra.
Megmutatjuk, hogy $B=C$.

Az $AB$ mátrix önadjungált és idempotens, mert
\[
  (AB)^*=AB, \qquad (AB)^2=ABAB=(ABA)B=AB.
\]
Továbbá $\Kep(AB)=\Kep(A)$, hiszen $\Kep(AB)\subseteq\Kep(A)$, míg
$ABA=A$ miatt $\Kep(A)\subseteq\Kep(AB)$. Ezért $AB$ az $\Kep(A)$-ra vett
ortogonális projekció. Ugyanez igaz $AC$-re is, tehát az ortogonális projekció
egyértelműsége miatt
\[
  AB=AC=:P.
\]
Hasonlóan $BA$ és $CA$ önadjungált idempotensek. Ráadásul
\[
  \Ker(BA)=\Ker(A),
\]
mert $Ax=0$ esetén nyilván $BAx=0$, míg $BAx=0$ esetén
$Ax=ABAx=0$. Így $BA$ és $CA$ ugyanarra az $\Ker(A)^\perp=\Kep(A^*)$ altérre
vetítenek ortogonálisan, tehát
\[
  BA=CA=:Q.
\]
Most bontsunk fel egy tetszőleges $y\in\C^m$ vektort az
$\C^m=\Kep(A)\oplus\Kep(A)^\perp$ ortogonális felbontás szerint: $y=y_R+y_N$.
Mivel $P=AB=AC$ az $\Kep(A)$-ra vett ortogonális projekció, $Py_N=0$, és
\[
  By_N=BAB y_N=B(AB)y_N=0,
  \qquad
  Cy_N=CAC y_N=C(AC)y_N=0.
\]
Ha pedig $y_R\in\Kep(A)$, akkor $y_R=Ax$ valamely $x$-re, és
\[
  By_R=BAx=Qx=CAx=Cy_R.
\]
Tehát $By=Cy$ minden $y$-ra, vagyis $B=C$. \qed
\end{proof}

\subsection*{A pszeudoinverz előállítása}

\begin{mdframed}[backgroundcolor=thmbg,linewidth=0pt,innertopmargin=8pt,innerbottommargin=8pt]
\begin{tetel}[Diagonális mátrix pszeudoinverze]
\label{tetel:diagpseudo}
Legyen $D \in \C^{m\times n}$ diagonális ($d_{ii}$-kkel). Ekkor $D^+ \in \C^{n\times m}$
szintén diagonális:
\[
  d_{ii}^+ = \begin{cases} 1/d_{ii}, & d_{ii} \ne 0, \\ 0, & d_{ii} = 0. \end{cases}
\]
\end{tetel}
\end{mdframed}

\begin{proof}
Ellenőrizzük az 1--4 feltételeket.
\begin{enumerate}[noitemsep,topsep=2pt]
  \item $(DD^+)_{ii} = d_{ii}\cdot d_{ii}^+ \in \{0, 1\}$; valós, tehát $(DD^+)^* = DD^+$. $\checkmark$
  \item Hasonlóan $(D^+D)^* = D^+D$. $\checkmark$
  \item $(DD^+D)_{ii} = d_{ii}\cdot d_{ii}^+\cdot d_{ii} = d_{ii}$. $\checkmark$
  \item $(D^+DD^+)_{ii} = d_{ii}^+\cdot d_{ii}\cdot d_{ii}^+ = d_{ii}^+$. $\checkmark$ \qed
\end{enumerate}
\end{proof}

\begin{mdframed}[backgroundcolor=thmbg,linewidth=0pt,innertopmargin=8pt,innerbottommargin=8pt]
\begin{tetel}[Az általánosított inverz előállítása SVD-vel]
\label{tetel:pseudosvd}
Legyen $A = UDV^*$ az $A$ szinguláris felbontása. Ekkor:
\[
  A^+ = V D^+ U^*,
\]
ahol $D^+$ a Tétel~\ref{tetel:diagpseudo} szerinti pszeudoinverz.
\end{tetel}
\end{mdframed}

\begin{proof}
Legyen $B = VD^+U^*$. Igazoljuk az 1--4 feltételeket:
\begin{enumerate}[noitemsep,topsep=2pt]
  \item $AB = UDV^*\cdot VD^+U^* = U(DD^+)U^*$; ez önadjungált, mert $DD^+$ önadjungált
        (Tétel~\ref{tetel:diagpseudo}).
  \item $BA = VD^+U^*\cdot UDV^* = V(D^+D)V^*$; szintén önadjungált.
  \item $ABA = UDV^*VD^+U^*UDV^* = U(DD^+D)V^* = UDV^* = A$.
  \item $BAB = VD^+U^*UDV^*VD^+U^* = V(D^+DD^+)U^* = VD^+U^* = B$. $\checkmark$
\end{enumerate}
\qed
\end{proof}

\subsection*{Speciális esetek}

\begin{mdframed}[backgroundcolor=thmbg,linewidth=0pt,innertopmargin=8pt,innerbottommargin=8pt]
\begin{tetel}[Túlhatározott, teljes rangú eset]
\label{tetel:tulhatarozottPS}
Legyen $A \in \C^{m\times n}$, $m > n$ és $\rang(A) = n$ (teljes oszloprangú). Ekkor:
\[
  A^+ = (A^*A)^{-1}A^*.
\]
\end{tetel}
\end{mdframed}

\begin{proof}
Használjuk a redukált SVD-t: $A=U_1\Sigma V^*$, ahol
$U_1\in\C^{m\times n}$ ortonormált oszlopú mátrix, $V\in\C^{n\times n}$ unitér,
és $\Sigma=\diag(\sigma_1,\ldots,\sigma_n)$ invertálható. Ekkor
\[
  A^+ = V\Sigma^{-1}U_1^*.
\]
Másrészt
\[
  A^*A = V\Sigma U_1^*U_1\Sigma V^* = V\Sigma^2V^*;
\]
így
\[
  (A^*A)^{-1}A^*
  = V\Sigma^{-2}V^*\,V\Sigma U_1^*
  = V\Sigma^{-1}U_1^* = A^+. \qed
\]
\end{proof}

\begin{mdframed}[backgroundcolor=thmbg,linewidth=0pt,innertopmargin=8pt,innerbottommargin=8pt]
\begin{tetel}[Alulhatározott, teljes rangú eset]
\label{tetel:alulhatarozottPS}
Legyen $A \in \C^{m\times n}$, $m < n$ és $\rang(A) = m$ (teljes sorrangú). Ekkor:
\[
  A^+ = A^*(AA^*)^{-1}.
\]
\end{tetel}
\end{mdframed}

\begin{proof}
Használjuk a redukált SVD-t: $A=U\Sigma V_1^*$, ahol
$U\in\C^{m\times m}$ unitér, $V_1\in\C^{n\times m}$ ortonormált oszlopú mátrix,
és $\Sigma=\diag(\sigma_1,\ldots,\sigma_m)$ invertálható. Ekkor
\[
  A^+ = V_1\Sigma^{-1}U^*.
\]
Másrészt
\[
  AA^* = U\Sigma V_1^*V_1\Sigma U^* = U\Sigma^2U^*;
\]
így
\[
  A^*(AA^*)^{-1}
  = V_1\Sigma U^*\,U\Sigma^{-2}U^*
  = V_1\Sigma^{-1}U^* = A^+. \qed
\]
\end{proof}

% ==============================================================
\section{Az általánosított inverz approximációs tulajdonsága}
\label{sec:approx}

\begin{mdframed}[backgroundcolor=defbg,linewidth=0pt,innertopmargin=8pt,innerbottommargin=8pt]
\begin{definicio}[Általánosított megoldás]
Az $Ax = b$ LER általánosított megoldása:
\[
  x^+ := A^+ b.
\]
\end{definicio}
\end{mdframed}

\subsection*{Következmények a speciális esetekben}

\begin{itemize}
  \item \textbf{Túlhatározott, teljes rangú eset} ($m > n$, $\rang(A) = n$):
    \[
      x^+ = (A^*A)^{-1}A^*b \quad \Longleftrightarrow \quad A^*A\,x^+ = A^*b.
    \]
    Az $A^*Ax = A^*b$ egyenletrendszert \emph{Gauss-féle normálegyenleteknek} nevezzük.
    $x^+$ a \emph{legkisebb négyzetek megoldása} (minimalizálja $\norm{Ax - b}_2$-t).

  \item \textbf{Alulhatározott, teljes rangú eset} ($m < n$, $\rang(A) = m$):
    \[
      x^+ = A^*(AA^*)^{-1}b.
    \]
    Ha $y := (AA^*)^{-1}b$, akkor $Ax^+ = AA^*y = b$, azaz $x^+$ valódi megoldás.
    Ráadásul $x^+$ az összes megoldás közül a \emph{minimális normájú}.
\end{itemize}

\begin{mdframed}[backgroundcolor=thmbg,linewidth=0pt,innertopmargin=8pt,innerbottommargin=8pt]
\begin{tetel}[Approximációs tulajdonság]
\label{tetel:approx}
Legyenek $A \in \C^{m\times n}$ és $b \in \C^m$ adottak, $x^+ = A^+b$.
\begin{enumerate}[noitemsep,topsep=2pt]
  \item $\norm{Ax - b}_2 \ge \norm{Ax^+ - b}_2 \quad \forall\,x \in \C^n$
        (azaz $x^+$ minimalizálja a maradékot).
  \item Ha $H := \{x \in \C^n : \norm{Ax - b}_2 = \norm{Ax^+ - b}_2\}$ az összes
        minimalizáló halmaza, akkor
        $\norm{x^+}_2 < \norm{x}_2 \quad \forall\,x \in H,\; x \ne x^+$
        (azaz $x^+$ minimális normájú a minimalizálók között).
\end{enumerate}
\end{tetel}
\end{mdframed}

\begin{proof}
\textbf{1. állítás.}
A Moore--Penrose-feltételekből következik, hogy $P:=AA^+$ önadjungált idempotens,
tehát ortogonális projekció. Mivel $AA^+A=A$, ez a projekció éppen az $\Kep(A)$ altérre
vetít. Így
\[
  Ax^+ = AA^+b = Pb
\]
a $b$ vektor ortogonális vetülete $\Kep(A)$-ra, ezért $b-Ax^+\perp\Kep(A)$.
Minden $x\in\C^n$ esetén $A(x-x^+)\in\Kep(A)$, tehát Pitagorasz tétele szerint
\[
  \norm{Ax-b}_2^2
  = \norm{A(x-x^+) + (Ax^+-b)}_2^2
  = \norm{A(x-x^+)}_2^2 + \norm{Ax^+-b}_2^2
  \ge \norm{Ax^+-b}_2^2.
\]
Ez igazolja, hogy $x^+$ minimalizálja a maradék normáját.

\textbf{2. állítás.}
Az előző azonosságban egyenlőség pontosan akkor áll fenn, ha
$A(x-x^+)=0$, vagyis $x-x^+\in\Ker(A)$. Az SVD-alakból látszik, hogy
$x^+=A^+b\in\Kep(A^*)$, és az alaptétel szerint $\Ker(A)=\Kep(A^*)^\perp$.
Ezért $x^+\perp(x-x^+)$, és ha $x\ne x^+$, akkor
\[
  \norm{x}_2^2
  = \norm{x^+ + (x-x^+)}_2^2
  = \norm{x^+}_2^2 + \norm{x-x^+}_2^2
  > \norm{x^+}_2^2.
\]
Tehát $x^+$ a minimális maradékú megoldások között egyértelműen a legkisebb normájú. \qed
\end{proof}

\begin{megjegyzes}
\leavevmode
\begin{enumerate}[noitemsep,topsep=2pt]
  \item Ha $Ax = b$ megoldható, akkor $x^+ = A^+b$ az origóhoz legközelebb lévő megoldás.
  \item Ha $Ax = b$ nem megoldható (például túlhatározott rendszer), akkor $x^+$ a
        \emph{reziduális norma} szerint legjobb közelítő megoldás, és ha több ilyen létezik,
        közülük a legkisebb normájú.
  \item A fenti két esetben (túl- és alulhatározott teljes rangú) a szinguláris felbontás
        nem szükséges: elegendő megoldani a normálegyenleteket illetve a minimál-norma rendszert.
\end{enumerate}
\end{megjegyzes}

% ==============================================================
\section{Legkisebb négyzetek módszere}
\label{sec:lnm}

\subsection*{Az alapfeladat}

\begin{mdframed}[backgroundcolor=defbg,linewidth=0pt,innertopmargin=8pt,innerbottommargin=8pt]
\begin{definicio}[A legkisebb négyzetek módszere]
Adottak az $x_1, \ldots, x_N \in [a,b]$ különböző alappontok és az
$y_1, \ldots, y_N \in \R$ függvény-, illetve mérési értékek.
Keresünk $p_n \in P_n$ polinomot ($n+1 \le N$, általában $N \gg n$), amelyre
\[
  \sum_{i=1}^N \bigl(y_i - p_n(x_i)\bigr)^2 \quad \text{minimális.}
\]
A $p_n$ polinomot \emph{négyzetesen legjobban közelítő polinomnak} nevezzük.
\end{definicio}
\end{mdframed}

A feladat tehát nem az, hogy $p_n(x_i) = y_i$ pontosan teljesüljön (ez csak $n = N-1$ esetén
érhető el), hanem az eltérések négyzetösszegét minimalizáljuk. Ez a statisztikában
\emph{lineáris regressziónak} felel meg.

\subsection*{Mátrix alak}

Írjuk fel $p_n(x) = a_n x^n + \cdots + a_1 x + a_0$ alakban. Jelöljük:
\[
  A :=
  \begin{pmatrix}
    1 & x_1 & x_1^2 & \cdots & x_1^n \\
    1 & x_2 & x_2^2 & \cdots & x_2^n \\
    \vdots & \vdots & \vdots & & \vdots \\
    1 & x_N & x_N^2 & \cdots & x_N^n
  \end{pmatrix}_{N\times(n+1)},
  \quad
  a :=
  \begin{pmatrix} a_0 \\ a_1 \\ \vdots \\ a_n \end{pmatrix}_{(n+1)\times 1},
  \quad
  y :=
  \begin{pmatrix} y_1 \\ y_2 \\ \vdots \\ y_N \end{pmatrix}_{N\times 1}.
\]
Ekkor a minimalizálandó mennyiség:
\[
  \norm{Aa - y}_2^2 = \sum_{i=1}^N (p_n(x_i) - y_i)^2.
\]
Ez pontosan a Tétel~\ref{tetel:approx} kerete: az $Aa = y$ LER általában
\emph{túlhatározott és nem megoldható} (ha $N > n+1$), de a Tétel szerint az
általánosított megoldás $a^+ = A^+ y$ minimalizálja a négyzethibát.

\subsection*{Gauss-féle normálegyenletek}

Ha $A$ teljes oszloprangú (rang$(A)=n+1$, ami teljesül, ha $N\ge n+1$ és az $x_i$ alappontok különbözők),
akkor Tétel~\ref{tetel:tulhatarozottPS} alapján:
\[
  a^+ = (A^T A)^{-1} A^T y \quad \Longleftrightarrow \quad A^T A\,a^+ = A^T y.
\]
A szimmetrikus, $(n+1)\times(n+1)$-es LER kifejtve:
\[
  \begin{pmatrix}
    N & \sum x_i & \cdots & \sum x_i^n \\
    \sum x_i & \sum x_i^2 & \cdots & \sum x_i^{n+1} \\
    \vdots & \vdots & \ddots & \vdots \\
    \sum x_i^n & \sum x_i^{n+1} & \cdots & \sum x_i^{2n}
  \end{pmatrix}
  \begin{pmatrix} a_0 \\ a_1 \\ \vdots \\ a_n \end{pmatrix}
  =
  \begin{pmatrix} \sum y_i \\ \sum x_i y_i \\ \vdots \\ \sum x_i^n y_i \end{pmatrix}.
\]

\begin{megjegyzes}
\leavevmode
\begin{enumerate}[noitemsep,topsep=2pt]
  \item Ha $A$ teljes rangú, $A^T A$ mindig szimmetrikus és pozitív definit, tehát invertálható.
  \item \textbf{Lineáris regresszió ($n=1$):}
    A normálegyenletek:
    \[
      \begin{pmatrix} N & \sum x_i \\ \sum x_i & \sum x_i^2 \end{pmatrix}
      \begin{pmatrix} a_0 \\ a_1 \end{pmatrix}
      = \begin{pmatrix} \sum y_i \\ \sum x_i y_i \end{pmatrix}.
    \]
    Ha $\sum x_i = 0$ (például szimmetrikus alappontok), diagonális a rendszer:
    $a_0 = \tfrac{1}{N}\sum y_i$, $a_1 = \tfrac{\sum x_i y_i}{\sum x_i^2}$.
    Általánosan:
    \[
      a_1 = \frac{N\sum x_i y_i - (\sum x_i)(\sum y_i)}{N\sum x_i^2 - (\sum x_i)^2},
      \qquad
      a_0 = \frac{\sum y_i - a_1 \sum x_i}{N}.
    \]
  \item A négyzetesen legjobban közelítő $p_1$ egyenes mindig átmegy az adatok
        \emph{átlagán} $(\bar{x}, \bar{y})$: ez éppen a normálegyenlet első sorából következik.
  \item A közgazdászok és statisztikusok ezt \emph{regressziós egyenesnek} hívják,
        és standardeszközként használják (statisztika, ökonometria).
\end{enumerate}
\end{megjegyzes}

\begin{figure}[h!]
\centering
\includegraphics[width=0.52\textwidth]{lin_regresszio.pdf}
\caption{Lineáris regresszió öt adatpontra. Az egyenes a legkisebb négyzetek módszerével
  illesztett $p_1(x) = a_1 x + a_0$ polinom; a szaggatott vonalak a reziduálisokat jelölik.}
\label{fig:linreg}
\end{figure}

\begin{mdframed}[backgroundcolor=exbg,linewidth=0pt,innertopmargin=8pt,innerbottommargin=8pt]
\textbf{Számpélda: lineáris regresszió.}
Adottak a következő mérési adatok:
\[
\begin{array}{c|ccccc}
x_i & 1 & 2 & 3 & 4 & 5 \\\hline
y_i & 2{,}1 & 3{,}8 & 6{,}2 & 7{,}9 & 9{,}5
\end{array}
\]
Számítsuk ki a legjobb lineáris közelítőt ($n=1$)!

$N = 5$, $\sum x_i = 15$, $\sum y_i = 29{,}5$, $\sum x_i^2 = 55$, $\sum x_i y_i = 100{,}3$.
\[
  a_1 = \frac{5 \cdot 100{,}3 - 15 \cdot 29{,}5}{5 \cdot 55 - 225}
  = \frac{501{,}5 - 442{,}5}{275 - 225} = \frac{59{,}0}{50} = 1{,}180.
\]
\[
  a_0 = \frac{29{,}5 - 1{,}180 \cdot 15}{5} = \frac{29{,}5 - 17{,}7}{5} = \frac{11{,}8}{5} = 2{,}360.
\]
Így $p_1(x) \approx 1{,}180\,x + 2{,}360$.

\textbf{Ellenőrzés:} A reziduális norma:
$\norm{Aa^+ - y}_2^2 = \sum (p_1(x_i)-y_i)^2 = \cdots \approx 0{,}022$ (számítható).

Az adatok átlaga: $\bar{x} = 3$, $\bar{y} = 5{,}9$; valóban $p_1(3) = 3{,}54 + 2{,}36 = 5{,}9$. $\checkmark$
\end{mdframed}

% ==============================================================
\section*{Összefoglalás}

\subsection*{Numerikus integrálás}

\begin{center}
\renewcommand{\arraystretch}{1.5}
\begin{tabular}{l|c|c|c}
\toprule
\textbf{Formula} & \textbf{Típus} & \textbf{Hiba} & \textbf{Pontossági fok}\\
\midrule
Érintő $E(f)$ & Ny(0) & $\dfrac{(b-a)^3}{24}f''(\eta)$ & $P_1$ \\
Trapéz $T(f)$ & Z(1) & $-\dfrac{(b-a)^3}{12}f''(\eta)$ & $P_1$ \\
Simpson $S(f)$ & Z(2) & $-\dfrac{(b-a)^5}{2880}f^{(4)}(\eta)$ & $P_3$ \\[4pt]
\midrule
Trapéz összetett $T_m(f)$ & — & $-\dfrac{(b-a)^3}{12m^2}f''(\eta)$ & $O(m^{-2})$ \\
Simpson összetett $S_m(f)$ & — & $-\dfrac{(b-a)^5}{180m^4}f^{(4)}(\eta)$ & $O(m^{-4})$ \\
\bottomrule
\end{tabular}
\end{center}

\subsection*{Szinguláris felbontás és általánosított inverz}

\begin{center}
\renewcommand{\arraystretch}{1.4}
\begin{tabular}{l|l}
\toprule
\textbf{Fogalom} & \textbf{Tartalom} \\
\midrule
SVD & $A = UDV^*$, $D = \diag(\sigma_1,\ldots,\sigma_r,0,\ldots)$, $U,V$ unitér \\
Szinguláris értékek & $\sigma_i = \sqrt{\lambda_i(A^*A)} > 0$ \\
Eckart--Young & $A_k = \sum_{i=1}^k \sigma_i u_i v_i^*$ a legjobb rang-$k$ közelítés \\
Moore--Penrose & $A^+ = VD^+U^*$ (4 ax.\ egyértelműen meghatározza) \\
Túlhatározott & $A^+ = (A^TA)^{-1}A^T$, normálegyenletek: $A^TAx = A^Ty$ \\
Alulhatározott & $A^+ = A^T(AA^T)^{-1}$, minimális norma \\
LNM & $\min \norm{Aa-y}_2^2$: $a^+ = (A^TA)^{-1}A^Ty$ (regresszió) \\
\bottomrule
\end{tabular}
\end{center}

\end{document}
